\documentclass[11pt,a4paper]{article}
\usepackage[T1]{fontenc}
\usepackage{lmodern}
\usepackage[margin=27mm,headheight=14pt]{geometry}
\usepackage{amsmath,amssymb,amsthm,mathtools,bm,mathrsfs}
\usepackage{microtype,booktabs,longtable,array,enumitem}
\usepackage{xcolor,hyperref,fancyhdr}
\usepackage{listings}
\hypersetup{colorlinks=true,linkcolor=blue!45!black,citecolor=blue!45!black,urlcolor=blue!45!black,pdftitle={Uniform q-Multinomial Transseries and Certified Inversion},pdfauthor={OpenAI, research draft prepared for Vladimir Reshetnikov}}
\pagestyle{fancy}\fancyhf{}
\fancyhead[L]{\small Uniform $q$-multinomial transseries}
\fancyhead[R]{\small Research draft}
\fancyfoot[C]{\thepage}
\setlength{\parskip}{3pt}
\setlength{\emergencystretch}{2em}
\numberwithin{equation}{section}
\newtheorem{theorem}{Theorem}[section]
\newtheorem{proposition}[theorem]{Proposition}
\newtheorem{lemma}[theorem]{Lemma}
\newtheorem{corollary}[theorem]{Corollary}
\theoremstyle{definition}\newtheorem{definition}[theorem]{Definition}
\newtheorem{example}[theorem]{Example}
\theoremstyle{remark}\newtheorem{remark}[theorem]{Remark}
\DeclareMathOperator{\Li}{Li}
\DeclareMathOperator{\Res}{Res}
\DeclareMathOperator{\arccoth}{arccoth}
\newcommand{\e}{\mathrm e}
\newcommand{\dd}{\,\mathrm d}
\newcommand{\N}{\mathbb N}
\newcommand{\R}{\mathbb R}
\newcommand{\C}{\mathbb C}
\newcommand{\avec}{\boldsymbol a}
\newcommand{\amin}{a_*}
\newcommand{\mmin}{m_*}
\newcommand{\calA}{\mathcal A}
\newcommand{\calP}{\mathcal P}
\newcommand{\calB}{\mathscr B}
\newcommand{\calL}{\mathcal L}
\newcommand{\qmulti}[2]{\genfrac{[}{]}{0pt}{}{#1}{#2}_{q}}
\newcommand{\repo}{\texttt{ProveIt}}
\lstset{basicstyle=\ttfamily\small,breaklines=true,columns=fullflexible,frame=single}

\title{\textbf{Uniform $q$-Multinomial Transseries\ and Certified Inversion}\\[5pt]
\large Closing the $q\to1$ endpoint gap, locating Borel singularities,\\
and proving a sharp least-term exponential rate}
\author{Research article prepared for Vladimir Reshetnikov\\
\small OpenAI}
\date{29 September 2026}

\begin{document}
\maketitle
\begin{abstract}
The combinatorial-transseries companion in \repo\ derives an all-order
$q\to1$ expansion of central Gaussian coefficients, but its estimates require
$\tau=hx\ge\tau_0>0$. It explicitly leaves uniform matching to $\tau=0$
for a different error analysis. We supply that analysis and extend it to every
positive multinomial ray. A partial-fraction identity for the Bose kernel gives
an exact signed remainder. After reorganizing the coefficients in powers of
$1/x$, the resulting estimates hold simultaneously for every $h\ge0$, with no
restriction on $hx$. The error is smaller than the first omitted term and is
bounded by an explicit Bernoulli expression independent of $h$.

We derive the meromorphic Borel transform, prove its exact Laplace-summation
formula, and show that the nearest singularities in the $1/x$ Borel plane are
always $\pm2\pi i a_*$, where $a_*$ is the smallest ray parameter and $m_*$ below is its multiplicity. A least-term
choice gives the uniform bound $2r\e^{-2\pi a_*x}$, retaining the modular
Euler-product contribution separately. At $h=0$ the actual remainder is
asymptotic to $m_*(2\pi\sqrt{a_*x})^{-1}\e^{-2\pi a_*x}$; hence its exponential
rate is sharp for this truncation rule. At the resonant values $h=2\pi/m$,
the remainder instead has equivalent $m_*(2\pi m)^{-1}\e^{-2\pi a_*x}$;
therefore no factor tending to zero with $x$ can multiply a bound uniform
in all $h$. A uniform lower bound on the exact
forward derivative transports these estimates to explicit inverse enclosures.
Complete proofs, endpoint formulas, finite evaluation bounds, numerical
checks, and ten further research directions are included.

\medskip\noindent
\textbf{Status.} The endpoint-uniform theorem package is developed and proved
here. It resolves an explicitly recorded repository gap, not a named general
conjecture about transseries. The elementary kernels and the general resurgence
framework have substantial prior literature. Global publication-level novelty
has not been established. This is an unrefereed research draft, not a
Lean-verified development.
\end{abstract}
\clearpage
\begingroup
\small
\setlength{\parskip}{0pt}
\tableofcontents
\endgroup
\newpage

\section{The research target and the precise contribution}
\label{sec:target}

\subsection{What the repository already supplies}
The relevant source is \emph{Combinatorial Transseries and Their Inverses},
in the series and transseries group of the Fabius-function documents
\cite{repo-companion,repo-readme}. Its chapter on the singular $q\to1$
transition uses $q=\e^h$, $Q=\e^{-h}$ and the finite-size coordinate $\tau=hx$.
It contains a dilogarithmic dominant phase, Bernoulli--polylogarithm corrections,
an exact modular Euler-product factor, and formal all-order inversion formulas.
Its proposition labelled \texttt{t3:prop:double-A} assumes $\tau\ge\tau_0>0$.
The paragraph following \texttt{t3:eq:Slarge} explicitly identifies the missing
uniform analysis at $\tau=0$. The modular-factor discussion also carefully
avoids claiming a full beyond-all-orders description of the Bernoulli tail.

Those are the two starting points of this article. We do not claim the
repository's dilogarithmic phase or its inverse coefficient formula as new.
Nor do we infer analytic theorems from a module inventory: the present proofs
are independent of the status of the repository's formalization.

The inspected snapshot is commit
\begin{center}\small\texttt{04e06e032dff1966513bfba316966d52db3646b4}.\end{center}
The companion's inspected source blob is
\begin{center}\small\texttt{77e494fbd30587b61b36f0e0e010c85cae3c7fa7}.\end{center}
The source ranges used most directly are lines 4080--4280, including the labels
just specified. The provenance file in this package records the full path.

\subsection{Why another finite-order expansion is not enough}
Three limits must be distinguished. At fixed $h>0$, finite-size corrections
are exponentially small in $x$. In the interior crossover, $h\to0$ and
$hx\to\tau\in(0,\infty)$; those corrections contribute to the leading
$1/h$ phase. At the classical endpoint, $h\to0$ with $hx\to0$ and $x\to\infty$,
the coefficient functions themselves become singular in the old coordinate.
A remainder constant depending on a positive lower bound for $hx$ cannot
control the third regime.

The cure is not to discard the singular coefficients. It is to absorb their
powers of $hx$ into the $1/x$ scale and to estimate the exact remainder before
passing to any limit. Positivity of the multinomial weight will preserve its
sign. This produces a single expansion with a rigorous meaning throughout
all three regimes, even when $h$ depends on $x$ in an arbitrary way.

\subsection{Relation to established analysis}
The partial fractions for $\coth$ and the eta modular transformation are
classical \cite{dlmf-hyperbolic,dlmf-eta}. Binet's integral and the Stirling
remainder are classical as well \cite{dlmf-gamma-int,dlmf-gamma-asymp}.
The resurgence of Euler--Maclaurin expansions was studied by Costin and
Garoufalidis \cite{costin-garoufalidis}; Garoufalidis and Kashaev give an explicit
Borel-plane description and Stokes analysis of Faddeev's quantum dilogarithm
\cite{garoufalidis-kashaev}. In particular, locating a lattice of Borel poles
for a $q$-product is not, in itself, a new phenomenon.

Our specific contribution is the combination of a positive multinomial
remainder, endpoint-uniform normalization, an $h$-independent error envelope,
the precise smallest-parameter Borel geometry, sharp endpoint and resonant
remainder equivalents, and a uniform inverse certificate. The targeted literature check
establishes relevant precedents, not an exhaustive priority search. The
statements below are mathematical claims with proofs; the assertion that this
package resolves the stated repository gap is narrower than an assertion of
worldwide priority.

\section{The family and the endpoint-uniform expansion}
\label{sec:definitions}

Fix $r\ge2$ and positive real numbers $\avec=(a_1,\ldots,a_r)$. Write
\begin{equation}
 A=\sum_{j=1}^r a_j,\qquad d=r-1,\qquad
 \kappa=\frac{A^2-\sum_j a_j^2}{2}>0,
 \qquad \amin=\min_j a_j.
 \label{eq:parameters}
\end{equation}
The multiplicity of the smallest parameter is
$\mmin=\#\{j:a_j=\amin\}$. Also put
\begin{equation}
 H_{\avec}=A\log A-\sum_j a_j\log a_j>0.
 \label{eq:entropy}
\end{equation}
The positivity follows on writing $a_j=A p_j$ and using
$H_{\avec}=-A\sum_j p_j\log p_j$.

For $h>0$, let
\begin{align}
 \calP(h)&=\prod_{n\ge1}(1-\e^{-hn}),\label{eq:product}\\
 \calA_h(s)&=\sum_{m\ge1}\frac{\e^{-ms}}{m(\e^{hm}-1)},
 \qquad s>0.\label{eq:tail}
\end{align}
Both converge absolutely. Define the logarithmic interpolation
\begin{equation}
 F_h(x)=\kappa h x^2-d\log\calP(h)
      +\calA_h(Ahx)-\sum_j\calA_h(a_jhx),\qquad x>0.
 \label{eq:exact-F}
\end{equation}
At $h=0$ define
\begin{equation}
 F_0(x)=\log\Gamma(Ax+1)-\sum_j\log\Gamma(a_jx+1).
 \label{eq:gamma-F}
\end{equation}
We prove continuity at $h=0$ below, rather than assuming it.

If all $a_jx$ are nonnegative integers, exponentiating \eqref{eq:exact-F}
gives the ordinary Gaussian multinomial coefficient
\begin{equation}
 \exp F_h(x)=\qmulti{Ax}{a_1x,\ldots,a_rx},\qquad q=\e^h.
 \label{eq:qmulti}
\end{equation}
Indeed, if $P_n=\prod_{j=1}^n(1-\e^{-hj})$, then
$P_n=\calP(h)\exp\calA_h(hn)$; extracting the powers of $q$ from the
Gaussian factorials gives the quadratic factor in \eqref{eq:exact-F}.
For arbitrary positive parameters, \eqref{eq:exact-F} is our specified real
interpolation. Nothing in the limiting argument requires a finite field whose
order tends to one.

\subsection{The phase, amplitude and modular term}
For $\tau>0$ set
\begin{align}
 S_{\avec}(\tau)
 &=\kappa\tau^2+\frac{d\pi^2}{6}
   +\Li_2(\e^{-A\tau})-\sum_j\Li_2(\e^{-a_j\tau}),\label{eq:S}\\
 f(u)&=\frac{1-\e^{-u}}{u},\qquad f(0)=1,\label{eq:f}\\
 B_{\avec}(\tau)
 &=\frac12\left(\log f(A\tau)-\sum_j\log f(a_j\tau)\right),\label{eq:B}\\
 M(h)&=-\log\calP(4\pi^2/h),\qquad M(0)=0.\label{eq:M}
\end{align}
All logarithms in these real formulas are real. The leading phase satisfies
$S_{\avec}(0)=0$ and
\begin{equation}
 S'_{\avec}(0)=H_{\avec}.
 \label{eq:S-endpoint}
\end{equation}
The modular term must not be silently discarded: it is flat as $h\downarrow0$
but is a nonzero constant in the fixed-$h$, large-$x$ problem.

Our unperturbed expression is
\begin{align}
 \Phi_h(x)={}&\frac{S_{\avec}(hx)}{h}
 -\frac d2\log(2\pi x)+\frac12\log\frac{A}{\prod_j a_j}
 +B_{\avec}(hx)-\frac{dh}{24}+dM(h),\quad h>0,\label{eq:core}\\
 \Phi_0(x)={}&H_{\avec}x-\frac d2\log(2\pi x)
              +\frac12\log\frac{A}{\prod_j a_j}.\label{eq:core-zero}
\end{align}
For $k\ge1$, the normalized coefficients are
\begin{equation}
 D_k(\tau)=\frac{B_{2k}}{(2k)!}\tau^{2k-1}
 \left(\Li_{2-2k}(\e^{-A\tau})
            -\sum_j\Li_{2-2k}(\e^{-a_j\tau})\right).
 \label{eq:D}
\end{equation}
Here $B_{2k}$ denotes a Bernoulli number, not the amplitude $B_{\avec}$.
At the endpoint the definition is continued by
\begin{equation}
 D_k(0)=\frac{B_{2k}}{2k(2k-1)}
       \left(A^{1-2k}-\sum_j a_j^{1-2k}\right).
 \label{eq:Dzero}
\end{equation}
Define
\begin{equation}
 T_K(h,x)=\Phi_h(x)+\sum_{k=1}^K\frac{D_k(hx)}{x^{2k-1}},
 \qquad R_K(h,x)=F_h(x)-T_K(h,x).
 \label{eq:TK}
\end{equation}
The empty sum defines $T_0=\Phi_h$.

\begin{theorem}[Endpoint-uniform signed expansion]
\label{thm:uniform}
For every positive ray $\avec$, every integer $K\ge0$, every $h\ge0$ and every
$x>0$, the expansion \eqref{eq:TK} obeys
\begin{equation}
 0<(-1)^{K+1}R_K(h,x)
   <\frac{|D_{K+1}(hx)|}{x^{2K+1}}
   \le E_K(x),
 \label{eq:main-bound}
\end{equation}
where
\begin{equation}
 E_K(x)=\frac{|B_{2K+2}|}{2K+2}\,
           x^{-(2K+1)}\sum_{j=1}^r a_j^{-(2K+1)}.
 \label{eq:EK}
\end{equation}
The coefficient functions are smooth at $\tau=0$, bounded on $[0,\infty)$,
and tend to zero as $\tau\to\infty$. The normalized core tends to
\eqref{eq:core-zero} as $h\downarrow0$ at fixed $x$.
Consequently the expansion is uniform along every choice $h=h(x)\ge0$ as
$x\to\infty$. In addition,
\begin{equation}
 |\partial_xR_K(h,x)|\le\frac{2K+2}{x}E_K(x),\qquad
 (-1)^K\partial_xR_K(h,x)>0.
 \label{eq:derivative-error}
\end{equation}
\end{theorem}

The theorem is a real-axis result. It makes no assertion of uniformity in a
complex sector whose Borel poles pinch the integration contour. That distinct
problem is discussed in Section~\ref{sec:questions}.

\section{The signed kernel and the proof of uniformity}
\label{sec:proof-uniform}

\subsection{An exact remainder before asymptotic expansion}
The classical partial fraction identity for $\coth$ \cite{dlmf-hyperbolic}
can be written, for $v>0$, as
\begin{equation}
 \frac1{\e^v-1}=\frac1v-\frac12
              +2v\sum_{\ell\ge1}\frac1{(2\pi\ell)^2+v^2}.
 \label{eq:bose-pf}
\end{equation}
Let $b_\ell=2\pi\ell$. A finite geometric expansion in each denominator gives
\begin{align}
 \frac1{\e^v-1}
 &=\frac1v-\frac12+
   \sum_{k=1}^K\frac{B_{2k}}{(2k)!}v^{2k-1}+r_K(v),\label{eq:kernel-exp}\\
 r_K(v)&=2(-1)^K v^{2K+1}
        \sum_{\ell\ge1}\frac{b_\ell^{-2K}}{b_\ell^2+v^2}.
 \label{eq:kernel-remainder}
\end{align}
The coefficient identity follows from
\begin{equation}
 \frac{B_{2k}}{(2k)!}
   =\frac{2(-1)^{k-1}\zeta(2k)}{(2\pi)^{2k}}.
 \label{eq:bernoulli-zeta}
\end{equation}
All series in this calculation converge absolutely for a fixed positive $v$.
Since every denominator in \eqref{eq:kernel-remainder} is positive,
\begin{equation}
 0<(-1)^K r_K(v)
 <\frac{|B_{2K+2}|}{(2K+2)!}v^{2K+1}.
 \label{eq:kernel-signed}
\end{equation}
No limiting interchange is needed to obtain this strict inequality.

\subsection{The multinomial weight}
The crucial positive function is
\begin{equation}
 W_{\avec}(s)=\sum_j\e^{-a_js}-\e^{-As}>0,
 \qquad s\ge0.
 \label{eq:W}
\end{equation}
At zero its value is $r-1$. For $s>0$, every $a_j<A$, so positivity is immediate.
Moreover,
\begin{equation}
 -W'_{\avec}(s)=\sum_j a_j\e^{-a_js}-A\e^{-As}
              =\sum_j a_j(\e^{-a_js}-\e^{-As})>0,\qquad s>0.
 \label{eq:W-derivative}
\end{equation}
Substitution of \eqref{eq:kernel-exp} into \eqref{eq:tail} now yields the exact
remainder identity
\begin{equation}
 R_K(h,x)=-\sum_{m\ge1}\frac{W_{\avec}(mhx)}{m}\,r_K(hm),
 \qquad h>0.
 \label{eq:Rexact}
\end{equation}
It simultaneously proves the sign and explains why taking separate absolute
values of the individual factorial remainders would be an unnecessary loss.

For completeness, the non-remainder terms assemble as follows. The $1/v$
term gives the dilogarithms divided by $h$; the constant $-1/2$ gives
\begin{equation}
 J_{\avec}(\tau)=\frac12\left(\log(1-\e^{-A\tau})
                  -\sum_j\log(1-\e^{-a_j\tau})\right).
 \label{eq:J}
\end{equation}
The eta modular identity \cite{dlmf-eta} is
\begin{equation}
 \calP(h)=\sqrt{\frac{2\pi}{h}}
          \exp\left(-\frac{\pi^2}{6h}+\frac h{24}\right)
          \calP(4\pi^2/h).
 \label{eq:eta}
\end{equation}
Finally,
\begin{equation}
 J_{\avec}(\tau)=B_{\avec}(\tau)-\frac d2\log\tau
                       +\frac12\log\frac{A}{\prod_j a_j}.
 \label{eq:J-normalized}
\end{equation}
Equations \eqref{eq:eta} and \eqref{eq:J-normalized} give precisely the core
\eqref{eq:core}. The finite Bernoulli terms give \eqref{eq:D}.

\subsection{A uniform moment estimate}
\begin{lemma}
\label{lem:moment}
For every integer $p\ge0$ and every $s>0$,
\begin{equation}
 s^{p+1}\Li_{-p}(\e^{-s})\le(p+1)!.
 \label{eq:moment-bound}
\end{equation}
\end{lemma}
\begin{proof}
On the interval $[(m-1)s,ms]$ one has $\e^{-t}\ge\e^{-ms}$, and hence
\begin{align*}
 \int_{(m-1)s}^{ms}t^p\e^{-t}\dd t
 &\ge\frac{s^{p+1}\e^{-ms}}{p+1}
          \bigl(m^{p+1}-(m-1)^{p+1}\bigr)\\
 &\ge\frac{s^{p+1}}{p+1}m^p\e^{-ms}.
\end{align*}
The second inequality follows from
$(m-1)^{p+1}\le(m-1)m^p$. Summing the integrals gives $p!$, proving the claim.
\end{proof}

Apply \eqref{eq:kernel-signed} in \eqref{eq:Rexact}. The resulting first-omitted
term is
\begin{align}
 |R_K(h,x)|<&\frac{|B_{2K+2}|}{(2K+2)!}h^{2K+1}
 \left(\sum_j\Li_{-2K}(\e^{-a_jhx})
                    -\Li_{-2K}(\e^{-Ahx})\right)\notag\\
 =&\frac{|D_{K+1}(hx)|}{x^{2K+1}}.
 \label{eq:next-bound}
\end{align}
Discard the negative last polylogarithm and use Lemma~\ref{lem:moment} at
$s=a_jhx$. The result is exactly \eqref{eq:EK}, with no dependence on $h$.
Differentiating \eqref{eq:Rexact} is justified by exponential convergence.
Equation \eqref{eq:W-derivative} gives the derivative's sign. The same argument
with $p=2K+1$ gives
\begin{equation}
 |\partial_xR_K|\le |B_{2K+2}|x^{-(2K+2)}
                               \sum_j a_j^{-(2K+1)},
\end{equation}
which is \eqref{eq:derivative-error}.

\subsection{Regularity of the coefficient functions}
The function $(\e^s-1)^{-1}$ is meromorphic near $s=0$, with principal part
$1/s$. Repeated differentiation shows
\begin{equation}
 \Li_{-p}(\e^{-s})=p!s^{-p-1}+\text{a function analytic at }s=0.
 \label{eq:li-endpoint}
\end{equation}
Multiplying by $\tau^{p+1}$ in \eqref{eq:D} removes the pole and gives
\eqref{eq:Dzero}. The coefficient functions thus extend analytically near zero.
At positive infinity they are a polynomial in $\tau$ times exponentially
small functions, so they tend to zero. Continuity on compact intervals then
proves boundedness.

Likewise $f$ is analytic and nonzero near zero, and the balanced relation
$A=\sum_j a_j$ cancels the linear term of its logarithm. Consequently
$B_{\avec}(\tau)=O(\tau^2)$. Differentiating \eqref{eq:S} gives
\begin{equation}
 S'_{\avec}(\tau)=2\kappa\tau+A\log(1-\e^{-A\tau})
                         -\sum_j a_j\log(1-\e^{-a_j\tau}).
 \label{eq:Sprime}
\end{equation}
Writing the logarithms as $\log(a_j\tau)+\log f(a_j\tau)$ proves
\eqref{eq:S-endpoint}. Also $M(h)\to0$, directly from its rapidly convergent
product representation. This establishes the core limit.

\subsection{The classical endpoint and continuity in \texorpdfstring{$h$}{h}}
Binet's integral for $z>0$ \cite{dlmf-gamma-int} reads
\begin{equation}
 \log\Gamma(z+1)=\left(z+\frac12\right)\log z-z
 +\frac12\log(2\pi)
 +\int_0^\infty\frac{\e^{-zt}}{t}
   \left(\frac1{\e^t-1}-\frac1t+\frac12\right)\dd t.
 \label{eq:Binet}
\end{equation}
Combining the gamma factors and using the same finite kernel expansion yields
\begin{equation}
 R_K(0,x)=2(-1)^{K+1}\sum_{\ell\ge1}b_\ell^{-2K}
 \int_0^\infty\frac{t^{2K}W_{\avec}(xt)}{b_\ell^2+t^2}\dd t.
 \label{eq:Rzero}
\end{equation}
Its sign is strict. Replacing its denominator by $b_\ell^2$ proves the strict
first-omitted-term estimate at $h=0$. The weaker bound $E_K(x)$ follows at once;
it also follows by taking the endpoint limit in the already normalized bound.
The derivative argument is identical.

There is a useful direct continuity check. Set $N=2K+1$ and $t=hm$ in
\eqref{eq:Rexact}. Its absolute value is
\begin{equation}
 2h\sum_{m\ge1}(hm)^{N-1}W_{\avec}(xhm)
         \sum_{\ell\ge1}\frac{b_\ell^{-(N-1)}}{b_\ell^2+(hm)^2}.
 \label{eq:R-riemann}
\end{equation}
This is a Riemann sum for the positive integral in \eqref{eq:Rzero}. The
integrand is continuous at zero and has exponential decay at infinity; the
$\ell$-sum is locally uniform and dominated by a convergent constant series.
A compact-interval Riemann-sum argument followed by an exponential tail bound
therefore gives convergence to \eqref{eq:Rzero}. Together with the finite
coefficient and core limits, this proves $F_h(x)\to F_0(x)$ for each $x>0$.
All assertions of Theorem~\ref{thm:uniform} are now established.

\section{Explicit specializations and the three limiting regimes}
\label{sec:examples}

\subsection{Central Gaussian coefficients}
For $\avec=(1,1)$, $A=2$, $\kappa=1$, $d=1$, and
\begin{align}
 S(\tau)&=\tau^2+\frac{\pi^2}{6}
                +\Li_2(\e^{-2\tau})-2\Li_2(\e^{-\tau}),\label{eq:Scentral}\\
 S'(\tau)&=2\log(1+\e^\tau),\label{eq:Scentral-prime}\\
 B(\tau)&=\frac12\log\left(\frac\tau2\coth\frac\tau2\right),\label{eq:Bcentral}\\
 \Phi_h(x)&=\frac{S(hx)}h-\frac12\log(\pi x)
                  +B(hx)-\frac h{24}+M(h).\label{eq:central-core}
\end{align}
Thus the repository phase is preserved exactly; the new feature is its
endpoint-uniform interpretation and error law.
The endpoint coefficients begin
\begin{equation}
 D_1(0)=-\frac18,\qquad D_2(0)=\frac1{192},\qquad
 D_3(0)=-\frac1{640},\qquad D_4(0)=\frac{17}{14336}.
 \label{eq:central-coeffs}
\end{equation}
Accordingly,
\begin{equation}
 \log\frac{\Gamma(2x+1)}{\Gamma(x+1)^2}
 \sim x\log4-\frac12\log(\pi x)-\frac1{8x}
             +\frac1{192x^3}-\frac1{640x^5}+\cdots.
\end{equation}
This classical expansion is not a novelty claim. It is the required endpoint
of the new uniform formula.

For every $h\ge0$, the first correction gives the rigorous enclosure
\begin{equation}
 0<F_h(x)-T_1(h,x)
   <\frac{|D_2(hx)|}{x^3}\le\frac1{60x^3}.
 \label{eq:first-simple-bound}
\end{equation}
The last constant is intentionally simple, not sharp: at $h=0$ the sharper
next-term bound is $1/(192x^3)$.

\subsection{A genuinely noncentral ray}
For $\avec=(1,2)$, the lattice values are $\qmulti{3x}{x,2x}$ and
\begin{equation}
 \kappa=2,\qquad H_{\avec}=3\log3-2\log2,\qquad
 D_1(0)=\frac1{12}\left(\frac13-1-\frac12\right)=-\frac7{72}.
\end{equation}
The smallest parameter has multiplicity one. This changes the sharp remainder
constant by a factor of two compared with the central ray, while leaving the
exponential rate $2\pi x$ unchanged. For $\avec=(1,1,1)$ the smallest parameter
has multiplicity three, so the sharp constant is three times the single-factor
one. These distinctions emerge directly from Borel residues, not from the
quadratic dominant phase.

\subsection{Classical, crossover, and fixed-\texorpdfstring{$q$}{q} limits}
At the classical endpoint,
\begin{equation}
 S_{\avec}(\tau)=H_{\avec}\tau+\frac\kappa2\tau^2
        +\frac{A^3-\sum_j a_j^3}{72}\tau^3+O(\tau^5).
 \label{eq:Ssmall}
\end{equation}
The term $\kappa\tau^2/(2h)$ explains the first deviation from the classical
multinomial phase. A regime such as $hx\to0$ does not imply that $hx^2\to0$;
retaining the complete phase avoids imposing this extra, unnecessary condition.

For $hx\to\tau\in(0,\infty)$, the terms in \eqref{eq:D} are bounded functions
of $\tau$ multiplying $x^{-1},x^{-3},\ldots$. This is the interior crossover
previously treated in the repository, now with explicit global constants.

For fixed $h>0$ and $x\to\infty$, the dilogarithms and all $D_k(hx)$ decay
exponentially. The core recombines by \eqref{eq:eta} into
$\kappa hx^2-d\log\calP(h)$ plus the appropriate finite-size terms. Thus no
spurious power correction is introduced into the fixed-$q$ problem. The
uniform bound is then often conservative, but remains valid.

\section{The exact Borel transform and its smallest action}
\label{sec:borel}

\subsection{Conventions and a single tail}
For a formal series without constant term use
\begin{equation}
 \calB\!\left(\sum_{n\ge0}c_nh^{n+1}\right)(\xi)
       =\sum_{n\ge0}\frac{c_n}{n!}\xi^n,
 \qquad (\calL\beta)(h)=\int_0^\infty\e^{-\xi/h}\beta(\xi)\dd\xi.
 \label{eq:Borel-convention}
\end{equation}
In particular there is no extra factor $1/h$ in the Laplace transform.
For $s>0$ define
\begin{equation}
 \beta_s(\xi)=\sum_{\ell\ge1}\frac1{b_\ell^2}
 \left(\frac1{\exp(s-i\xi/b_\ell)-1}
       +\frac1{\exp(s+i\xi/b_\ell)-1}\right).
 \label{eq:beta-single}
\end{equation}

\begin{theorem}[Meromorphic Borel summation]
\label{thm:borel}
The series \eqref{eq:beta-single} is locally uniformly convergent away from its
poles and is the Borel transform of
\begin{equation}
 \sum_{k\ge1}\frac{B_{2k}}{(2k)!}\Li_{2-2k}(\e^{-s})h^{2k-1}.
 \label{eq:formal-single}
\end{equation}
Its poles are simple and lie at
\begin{equation}
 \xi=4\pi^2\ell j\ \pm\ 2\pi i\ell s,
 \qquad \ell\ge1,\quad j\in\mathbb Z.
 \label{eq:single-poles}
\end{equation}
The residue at an upper-half-plane pole is $-i/(2\pi\ell)$ and at a
lower-half-plane pole is $+i/(2\pi\ell)$. For every $h,s>0$,
\begin{equation}
 \calA_h(s)=\frac{\Li_2(\e^{-s})}{h}
      +\frac12\log(1-\e^{-s})+(\calL\beta_s)(h).
 \label{eq:exact-borel}
\end{equation}
\end{theorem}
\begin{proof}
In a neighborhood of $\xi=0$, expand the two denominators in exponentials:
\begin{equation}
 \beta_s(\xi)=2\sum_{\ell\ge1}b_\ell^{-2}
                       \sum_{m\ge1}\e^{-ms}\cos(m\xi/b_\ell).
 \label{eq:beta-cos}
\end{equation}
The Taylor coefficient of $\xi^{2k-2}$ is
\[
 \frac{2(-1)^{k-1}}{(2k-2)!}
       \sum_{\ell\ge1}b_\ell^{-2k}
       \sum_{m\ge1}m^{2k-2}\e^{-ms},
\]
which is exactly the Borel coefficient in \eqref{eq:formal-single} by
\eqref{eq:bernoulli-zeta}. For a compact set of $\xi$ values, the terms with
large $\ell$ are uniformly $O(\ell^{-2})$. Only finitely many remaining terms
can have poles in that compact set. This proves meromorphic continuation and
local uniform convergence away from the poles. Solving the denominator equations
and differentiating them gives \eqref{eq:single-poles} and the stated residues.

On the positive real axis the absolute value of \eqref{eq:beta-cos} is at most
$2\zeta(2)(2\pi)^{-2}(\e^s-1)^{-1}$. Hence its Laplace integral is absolutely
convergent and summation may be interchanged with integration. Since
\[
 \int_0^\infty\e^{-\xi/h}\cos(m\xi/b_\ell)\dd\xi
              =\frac{h}{1+(hm/b_\ell)^2},
\]
the result is the partial-fraction correction in \eqref{eq:bose-pf}, inserted
into \eqref{eq:tail}. This proves \eqref{eq:exact-borel}.
\end{proof}

\subsection{The multinomial Borel plane}
At fixed $\tau>0$ put
\begin{equation}
 \beta_{\avec,\tau}(\xi)=\beta_{A\tau}(\xi)-\sum_j\beta_{a_j\tau}(\xi).
 \label{eq:beta-multi}
\end{equation}
The correct Borel transform for the normalized $1/x$ series is
\begin{equation}
 \calB_{\avec,\tau}(\eta)=\tau\beta_{\avec,\tau}(\tau\eta).
 \label{eq:scaled-Borel}
\end{equation}
The prefactor $\tau$ is essential: it accounts for replacing $h$ by $\tau/x$.
The exact summation identity becomes
\begin{equation}
 F_h(x)-\Phi_h(x)=\int_0^\infty\e^{-x\eta}
                              \calB_{\avec,hx}(\eta)\dd\eta.
 \label{eq:scaled-Laplace}
\end{equation}

\begin{theorem}[Universal nearest Borel singularities]
\label{thm:nearest}
For every $\tau>0$, the Taylor series of $\calB_{\avec,\tau}$ at zero has
radius of convergence exactly $2\pi\amin$. Its nearest singularities are
\begin{equation}
 \eta=\pm2\pi i\amin,
 \label{eq:nearest}
\end{equation}
with residues $+i\mmin/(2\pi)$ at the upper pole and $-i\mmin/(2\pi)$ at the
lower pole. They cannot cancel.
Other candidate poles have the form
\begin{equation}
 \eta=\frac{4\pi^2\ell j}{\tau}\ \pm\ 2\pi i\ell\lambda,
 \qquad \lambda\in\{A,a_1,\ldots,a_r\}.
 \label{eq:scaled-poles}
\end{equation}
At coincident poles the residues are added, and cancellation is possible.
As $\tau\downarrow0$, the meromorphic functions converge locally uniformly,
away from the limiting poles, to
\begin{equation}
 \calB_{\avec,0}(\eta)=\sum_{\ell\ge1}
 \left(\frac{2A}{b_\ell^2A^2+\eta^2}
           -\sum_j\frac{2a_j}{b_\ell^2a_j^2+\eta^2}\right).
 \label{eq:Borel-gamma}
\end{equation}
\end{theorem}
\begin{proof}
Equation \eqref{eq:scaled-poles} follows by scaling
\eqref{eq:single-poles}. Every such point has modulus at least
$2\pi\ell\lambda\ge2\pi\amin$. Equality is possible only when $j=0$,
$\ell=1$ and $\lambda=\amin$. The numerator parameter $A$ is strictly greater
than $\amin$. Thus exactly the $\mmin$ negative denominator contributions
occur there, and their residues add rather than cancel. This proves the exact
radius and the residues.

For the endpoint limit, each finite-$\ell$ term uses
\[
 \frac{\tau}{\exp\{\tau(\lambda\pm i\eta/b_\ell)\}-1}
       \longrightarrow\frac1{\lambda\pm i\eta/b_\ell}.
\]
On a fixed compact set and for sufficiently large $\ell$, the real part of
$\lambda\pm i\eta/b_\ell$ is at least $\lambda/2$. The corresponding terms
are uniformly bounded by a constant times $\ell^{-2}$, for all sufficiently
small positive $\tau$. Dominated convergence of the tail, together with
uniform convergence of finitely many analytic terms away from their limiting
poles, proves \eqref{eq:Borel-gamma}.
\end{proof}

This gives a precise geometric interpretation of the endpoint normalization.
The poles with $j\ne0$ move off to infinity as $\tau\to0$, while the nearest
imaginary poles remain fixed. The limiting kernel is exactly the Borel kernel
of the gamma-multinomial Stirling correction. Endpoint uniformity therefore
corresponds to a controlled meromorphic limit, not to removal of the leading
factorial divergence.

\subsection{Large-order coefficients and what is not being claimed}
For each fixed $\tau>0$, the residue pair in Theorem~\ref{thm:nearest} gives
\begin{equation}
 D_k(\tau)\sim
 \frac{2\mmin(-1)^k(2k-2)!}
      {(2\pi)^{2k}\amin^{2k-1}},\qquad k\to\infty.
 \label{eq:large-order}
\end{equation}
The same equivalent holds at $\tau=0$. To justify it directly, subtract the
two nearest principal parts from \eqref{eq:scaled-Borel}. For fixed $\tau$,
there is a strictly larger disk free of the remaining singularities; Cauchy
estimates make their coefficient contribution exponentially smaller in $k$.
Expanding the two principal parts gives \eqref{eq:large-order}.

The relative error implicit in \eqref{eq:large-order} is not asserted to be
uniform as $\tau\to\infty$, when other poles approach the same modulus. The
uniform truncation bound in Theorem~\ref{thm:uniform} does not depend on such
a coefficient equivalent and remains valid there.

The pole lattice determines residue contributions under an admissible
contour deformation: a positively oriented small contour around a pole
$\omega$ contributes
$2\pi i\,\Res(\beta_{\avec,\tau},\omega)\e^{-\omega/h}$.
This local statement is exact. A full global Stokes theorem would additionally
specify lateral contours, growth estimates and all collision cases. We do not
replace those requirements by a formal list of exponentials.

\section{Uniform exponential accuracy and a sharp endpoint equivalent}
\label{sec:optimal}

\begin{theorem}[An explicit least-term envelope]
\label{thm:exponential}
Let $z=\amin x\ge1$. Choose an odd integer $N$ nearest to $2\pi z$, breaking
an exact tie either way, and set $K=(N-1)/2$. Then for every $h\ge0$,
\begin{equation}
 |R_K(h,x)|\le E_K(x)
       <3r\sqrt z\,\e^{-2\pi z}.
 \label{eq:exponential}
\end{equation}
In particular, for the central Gaussian ray,
\begin{equation}
 |R_K(h,x)|<6\sqrt x\,\e^{-2\pi x},\qquad h\ge0,\quad x\ge1.
 \label{eq:central-exponential}
\end{equation}
\end{theorem}
\begin{proof}
Because $a_j\ge\amin$, \eqref{eq:EK} and \eqref{eq:bernoulli-zeta} imply
\begin{equation}
 E_K(x)\le\frac{2r\zeta(N+1)}{2\pi}\,
                  \frac{N!}{(2\pi z)^N}.
 \label{eq:factorial-bound}
\end{equation}
Write $u=2\pi z$ and $N=u+\delta$, $|\delta|\le1$. The elementary real Stirling
upper bound $N!\le\sqrt{2\pi N}(N/\e)^N\e^{1/(12N)}$ follows, for example,
from Binet's positive first remainder \cite{dlmf-gamma-asymp}. The function
$v\log(v/u)-v$ has value $-u$ and derivative zero at $v=u$, with second
derivative $1/v$. Hence
\[
 N\log(N/u)-N\le-u+\frac1{2(u-1)}.
\]
It follows that
\begin{equation}
 \frac{N!}{u^N}\le\sqrt{2\pi(u+1)}\,
        \exp\left(-u+\frac7{12(u-1)}\right).
 \label{eq:near-stirling}
\end{equation}
Here $u\ge2\pi$, $N\ge5$, and $\zeta(N+1)\le\zeta(6)<1.02$.
Dividing the combination of \eqref{eq:factorial-bound} and
\eqref{eq:near-stirling} by $r\sqrt z\e^{-2\pi z}$ bounds the ratio by
\[
 2(1.02)\sqrt{1+\frac1{2\pi}}
           \exp\left(\frac7{12(2\pi-1)}\right)<3.
\]
This proves the stated constant without an unspecified asymptotic threshold.
\end{proof}

\subsection{Removing the polynomial loss uniformly}
The preceding bound follows directly from the simple all-order expression
$E_K$. A sharper lattice-moment estimate removes its factor $\sqrt{x}$.

\begin{lemma}[A peak-sensitive lattice moment]
\label{lem:peak-moment}
For every integer $N\ge5$ and every $s>0$,
\begin{equation}
 s^N\Li_{1-N}(\e^{-s})\le 3N^N\e^{-N}.
 \label{eq:peak-moment}
\end{equation}
\end{lemma}
\begin{proof}
If $f$ is a nonnegative integrable unimodal function on $[0,\infty)$,
\begin{equation}
 s\sum_{m\ge1}f(ms)\le\int_0^\infty f(t)\dd t+s\sup f.
 \label{eq:unimodal-sum}
\end{equation}
Indeed, every positive superlevel set is an interval. The number of positive
lattice points of spacing $s$ in an interval is at most its length divided
by $s$, plus one. Integrating this inequality over the superlevel height
proves \eqref{eq:unimodal-sum} by the layer-cake formula.

Apply it to $f(t)=t^{N-1}\e^{-t}$. If $s\le N+1$, the left side of
\eqref{eq:peak-moment} is at most
\[
 \Gamma(N)+(N+1)(N-1)^{N-1}\e^{-(N-1)}.
\]
After division by $N^N\e^{-N}$, Stirling's upper bound makes the first term
at most $\sqrt{2\pi/N}\e^{1/(12N)}<1.15$. The second term is
\[
 \e(1+1/N)(1-1/N)^{N-1}
 \le(1+1/N)\e^{1/N}\le1.2\e^{1/5}<1.47.
\]
Their sum is less than 3. If $s\ge N+1$, successive terms of
$\sum m^{N-1}\e^{-ms}$ have ratio at most
$2^{N-1}\e^{-s}\le16\e^{-6}<1/25$. Thus
\[
 s^N\Li_{1-N}(\e^{-s})
 \le\frac{25}{24}s^N\e^{-s}
 \le\frac{25}{24}N^N\e^{-N}<3N^N\e^{-N}.
\]
The last inequality uses the maximum of $t^N\e^{-t}$ at $t=N$.
\end{proof}

\begin{theorem}[Uniform exponential bound without a polynomial loss]
\label{thm:constant-envelope}
With $z$, $N$ and $K$ chosen as in Theorem~\ref{thm:exponential},
\begin{equation}
 \boxed{\quad |R_K(h,x)|<2r\e^{-2\pi z},
           \qquad h\ge0,\quad z=\amin x\ge1.\quad}
 \label{eq:constant-envelope}
\end{equation}
\end{theorem}
\begin{proof}
Use the first-omitted-term inequality in \eqref{eq:main-bound} and discard
the negative numerator contribution in its positive weight. For $h>0$,
Lemma~\ref{lem:peak-moment} and \eqref{eq:bernoulli-zeta} give
\begin{align*}
 |R_K(h,x)|
 &<\frac{2\zeta(N+1)}{(2\pi)^{N+1}}
     h^N\sum_j\Li_{1-N}(\e^{-a_jhx})\\
 &\le\frac{3r\zeta(N+1)}{\pi}
                   \left(\frac{N}{2\pi z}\right)^N\e^{-N}.
\end{align*}
The same estimate at $h=0$ follows by continuity. Writing $u=2\pi z$ and
using $|N-u|\le1$ as before, the right side is at most
\[
 \frac{3r\zeta(6)}{\pi}
       \exp\left(-u+\frac1{2(u-1)}\right)<2r\e^{-u},
 \qquad u\ge2\pi.
\]
The numerical constant 2 is deliberately simple rather than best possible.
\end{proof}

The exponential here is a bound on optimally truncated remainders. The nearest
Borel singularities are imaginary, not positive real. It would be incorrect
to describe $\e^{-2\pi\amin x}$ as a positive-real Stokes sector solely from
this estimate. The following theorem gives the precise sharpness assertion.

\begin{theorem}[Sharp endpoint remainder]
\label{thm:sharp}
Fix a positive ray $\avec$. Use the same nearest-odd rule for $N$ and $K$ as
in Theorem~\ref{thm:exponential}. As $x\to\infty$,
\begin{equation}
 |R_K(0,x)|\sim\frac{\mmin}{2\pi\sqrt{\amin x}}
                         \e^{-2\pi\amin x}.
 \label{eq:sharp}
\end{equation}
Consequently, for this least-term truncation rule, no uniform estimate over
$h\ge0$ can replace the exponential rate $2\pi\amin$ by
$2\pi\amin+\varepsilon$, even allowing an arbitrary fixed polynomial
prefactor in $x$.
\end{theorem}
\begin{proof}
Let $b=2\pi$ and $z=\amin x$. In \eqref{eq:Rzero}, the terms with $\ell\ge2$
are exponentially smaller in $N$: their denominator is larger, and the
coefficient ratio to $\ell=1$ is at most $\ell^{-(N-1)}$. The sum of these
ratios is $O(2^{-N})$.

Within the $\ell=1$ integral, the terms with $a_j>\amin$, and the numerator
term with $A>\amin$, are likewise negligible. To see this quantitatively,
write
\begin{equation}
 I_N(y)=\int_0^\infty\frac{t^{N-1}\e^{-yt}}{b^2+t^2}\dd t
       =\frac{\Gamma(N)}{y^N}\,
            \mathbb E\frac1{b^2+T_y^2},
 \label{eq:gamma-expectation}
\end{equation}
where $T_y$ has gamma shape $N$ and rate $y$. If $y=z$, its mean is
$N/z\to b$ and its variance is $N/z^2\to0$. The bounded continuous integrand
therefore has expectation tending to $1/(2b^2)$. In particular this expectation
is bounded below by a positive constant for large $x$. For $y=a_jx>z$, its
expectation is bounded above by $1/b^2$, so
$I_N(a_jx)/I_N(z)=O((\amin/a_j)^N)$. The same applies to $A$.

Thus \eqref{eq:Rzero} is asymptotic in absolute value to
\begin{align}
 2\mmin b^{-(N-1)}I_N(z)
 &\sim\frac{\mmin\Gamma(N)}{b^{N+1}z^N}.\label{eq:sharp-middle}
\end{align}
Since $N=bz+O(1)$, Stirling's formula gives
\[
 \frac{\Gamma(N)}{(bz)^N}
     \sim\sqrt{\frac{2\pi}{bz}}\e^{-bz}.
\]
Substitute $b=2\pi$ in \eqref{eq:sharp-middle} to obtain \eqref{eq:sharp}.
The final assertion follows by evaluating any proposed stronger uniform bound
at $h=0$: an extra factor $\e^{-\varepsilon x}$ cannot dominate the nonzero
right-hand side of \eqref{eq:sharp} up to a fixed polynomial.
\end{proof}

\subsection{Discrete resonances obstruct an endpoint-sized uniform prefactor}
The $x^{-1/2}$ factor at $h=0$ comes from a continuous saddle integral.
At fixed $h>0$, the exact remainder is a lattice sum. A lattice point can
coincide with that saddle and remove the shrinking Gaussian-width factor.
This gives a second sharpness result.

\begin{theorem}[Resonant fixed-$h$ remainder]
\label{thm:resonance}
Fix an integer $m_0\ge1$ and set $h=2\pi/m_0$. For a fixed positive ray,
and the same nearest-odd rule as above,
\begin{equation}
 |R_K(2\pi/m_0,x)|\sim\frac{\mmin}{2\pi m_0}
                  \e^{-2\pi\amin x},\qquad x\to\infty.
 \label{eq:resonance}
\end{equation}
In particular, no estimate of the form
$C_{\avec}g(x)\e^{-2\pi\amin x}$ with $g(x)\to0$ can hold uniformly for
all $h\ge0$ for this truncation rule. Combined with
Theorem~\ref{thm:constant-envelope}, this proves that the uniform remainder
has precisely the order $\e^{-2\pi\amin x}$ up to positive ray-dependent
constant factors.
\end{theorem}
\begin{proof}
Write $b=2\pi$, $z=\amin x$ and $N=bz+\delta$, $|\delta|\le1$.
The absolute version of \eqref{eq:Rexact} is
\begin{equation}
 |R_K(h,x)|=2h\sum_{m\ge1}(mh)^{N-1}W_{\avec}(mhx)
       \sum_{\ell\ge1}\frac{(b\ell)^{-(N-1)}}{(b\ell)^2+(mh)^2}.
 \label{eq:discrete-remainder}
\end{equation}
The $\ell\ge2$ contributions are at most
$\sum_{\ell\ge2}\ell^{-(N-1)}=O(2^{-N})$ times the $\ell=1$
contribution, since every weight is positive.

First retain only the $\mmin\e^{-\amin mhx}$ part of $W_{\avec}$.
At $m=m_0$ the product $mh$ equals $b$ exactly, and its contribution for
$\ell=1$ is
\[
 \frac{\mmin h}{b^2}\e^{-bz}
       =\frac{\mmin}{b m_0}\e^{-bz}.
\]
For any other $m$, put $v=m/m_0$. The ratio of its contribution to the
$m_0$ contribution is
\begin{equation}
 \frac{2v^{\delta-1}}{1+v^2}
       \exp\{-bz(v-1-\log v)\}.
 \label{eq:resonance-ratio}
\end{equation}
The function $v-1-\log v$ is positive away from $v=1$. On this fixed lattice
its infimum for $m\ne m_0$ is positive; it grows at least linearly for large
$m$. The prefactor is bounded, and summing the exponentially decaying tail
shows that the sum of \eqref{eq:resonance-ratio} is $O(\e^{-c z})$ for some
$c>0$ depending on $m_0$.

It remains to bound the terms in $W_{\avec}$ with parameter
$\lambda>\amin$, including its negative numerator term. Replacing the
rational denominator in \eqref{eq:discrete-remainder} by $b^{-2}$ and
using Lemma~\ref{lem:peak-moment} bounds each such absolute contribution by
\[
 \frac{6}{b}\left(\frac{N}{b\lambda x}\right)^N\e^{-N}
       =O\left(\e^{-bz}(\amin/\lambda)^N\right).
\]
This is negligible relative to the displayed $m_0$ contribution. It proves
\eqref{eq:resonance}. The obstruction follows already from $m_0=1$.
The lower bound on the supremum over $h$ is supplied by that same value;
Theorem~\ref{thm:constant-envelope} supplies the upper bound.
\end{proof}

\begin{remark}
Both sharpness claims concern the specified Bernoulli truncation rule, not
all possible numerical algorithms or hyperasymptotic corrections. Subtracting
a resolved leading remainder can improve the error. The resonances occur at
fixed positive $h$, so they do not contradict the endpoint equivalent or a
possible improved estimate restricted to a sufficiently fine lattice
$h=h(x)\to0$. Such a restriction requires quantitative control of the lattice
spacing relative to the saddle width.
\end{remark}

\section{Uniform conditioning and inverse certificates}
\label{sec:inverse}

An asymptotic formula for a forward function is not automatically an equally
accurate inverse formula. A small derivative can amplify a residual. Here the
multinomial structure gives a derivative bound valid uniformly in $h$.

\begin{theorem}[Uniform lower bound on the inverse slope]
\label{thm:slope}
Set
\begin{equation}
 x_0=\amin^{-1},\qquad
 c_{\avec}=A\psi(Ax_0+1)-\sum_j a_j\psi(a_jx_0+1)>0,
 \label{eq:c}
\end{equation}
where $\psi$ is the digamma function. Then
\begin{equation}
 F_h'(x)\ge F_0'(x)\ge c_{\avec},
 \qquad h\ge0,\quad x\ge x_0.
 \label{eq:slope}
\end{equation}
For $\avec=(1,1)$ one has $x_0=1$ and $c_{\avec}=1$.
\end{theorem}
\begin{proof}
Differentiate \eqref{eq:exact-F} at fixed $h$ and put
\begin{equation}
 g_x(t)=\frac{\sum_j a_j\e^{-a_jxt}-A\e^{-Axt}}{\e^t-1}.
 \label{eq:g}
\end{equation}
Then
\begin{equation}
 F_h'(x)=2\kappa hx+h\sum_{m\ge1}g_x(hm),\qquad
 g_x(0)=2\kappa x.
 \label{eq:Fprime-sum}
\end{equation}
For $x\ge x_0$, $g_x$ is decreasing. Indeed it is the sum of the positive
functions
\begin{equation}
 a_j\e^{-(a_jx+1)t}
       \frac{1-\e^{-(A-a_j)xt}}{1-\e^{-t}}.
 \label{eq:g-positive}
\end{equation}
Here $(A-a_j)x\ge1$. For any $v\ge1$,
\[
 \frac{\dd}{\dd t}\log\frac{1-\e^{-vt}}{1-\e^{-t}}
      =\frac{v}{\e^{vt}-1}-\frac1{\e^t-1}\le0,
\]
because $v\mapsto v/(\e^{vt}-1)$ is decreasing. Thus every summand in
\eqref{eq:g-positive} is decreasing.

The lower Riemann sum inequality is
\[
 h\sum_{m\ge1}g_x(hm)\ge\int_0^\infty g_x(t)\dd t-hg_x(0).
\]
The subtracted term cancels the quadratic derivative in
\eqref{eq:Fprime-sum}. The remaining integral is $F_0'(x)$, by the elementary
integral representation of digamma differences. Hence $F_h'(x)\ge F_0'(x)$.

It remains to show that $F_0'$ is strictly increasing and positive. The
trigamma integral gives
\begin{equation}
 F_0''(x)=\int_0^\infty s\e^{-xs}
 \left(\frac1{\e^{s/A}-1}
             -\sum_j\frac1{\e^{s/a_j}-1}\right)\dd s.
 \label{eq:F0-second}
\end{equation}
For fixed $s>0$, the function
$u\mapsto [u(\e^{s/u}-1)]^{-1}$ is strictly increasing. Since $a_j<A$,
\[
 \sum_j\frac1{\e^{s/a_j}-1}
 <\sum_j\frac{a_j}{A}\frac1{\e^{s/A}-1}
 =\frac1{\e^{s/A}-1}.
\]
Therefore $F_0''(x)>0$. Also $F_0'(0)=0$, since $A=\sum_j a_j$.
This proves \eqref{eq:c} and \eqref{eq:slope}. In the central case,
$2\psi(3)-2\psi(2)=1$.
\end{proof}

For $h>0$, the numerator in \eqref{eq:g} is positive for every $x>0$.
Thus $F_h$ is increasing already below $x_0$; the threshold is used only for
the uniform numerical slope bound. Moreover $F_h(0)=0$ by continuous extension
of \eqref{eq:exact-F}, and $F_h(x)\to\infty$. The same statements hold for
$F_0$. Each has an unambiguous inverse on positive logarithmic targets.

\begin{theorem}[Residual-to-inverse transport]
\label{thm:certificate}
Let $L\ge F_h(x_0)$, and let $x_*$ be the unique solution of $F_h(x_*)=L$.
For any $s\ge x_0$, put
\begin{equation}
 \rho=T_K(h,s)-L,\qquad e=E_K(s).
\end{equation}
Then
\begin{equation}
 |x_*-s|\le\frac{|\rho|+e}{c_{\avec}}.
 \label{eq:general-certificate}
\end{equation}
When $\rho=0$, the sharper one-sided enclosures are
\begin{align}
 s\le x_*\le s+e/c_{\avec},&\qquad K\ \text{even},\label{eq:inverse-even}\\
 \max\{x_0,s-e/c_{\avec}\}\le x_*\le s,&\qquad K\ \text{odd}.
 \label{eq:inverse-odd}
\end{align}
One may replace $e$ throughout by the smaller next-term bound in
\eqref{eq:main-bound}. If $K$ is chosen by the least-term rule at $s$, then
\begin{equation}
 |x_*-s|\le\frac{|\rho|+2r\e^{-2\pi\amin s}}
                    {c_{\avec}}.
 \label{eq:inverse-exp}
\end{equation}
\end{theorem}
\begin{proof}
The exact residual is $F_h(s)-L=\rho+R_K(h,s)$. Theorem~\ref{thm:uniform}
bounds its magnitude by $|\rho|+e$. Both $s$ and $x_*$ lie in $[x_0,\infty)$,
so the mean value theorem and Theorem~\ref{thm:slope} give
\eqref{eq:general-certificate}. When $\rho=0$, the sign in
\eqref{eq:main-bound} determines which side of $s$ contains $x_*$. Finally use
Theorem~\ref{thm:constant-envelope}.
\end{proof}

\subsection{Existence and monotonicity of truncated inverse equations}
For an odd $K$, \eqref{eq:derivative-error} says $\partial_xR_K<0$.
Consequently
\begin{equation}
 \partial_xT_K(h,x)=F_h'(x)-\partial_xR_K(h,x)\ge c_{\avec}
 \quad(x\ge x_0).
\end{equation}
Also $T_K(h,x_0)<F_h(x_0)$ and $T_K(h,x)\to\infty$. Thus an odd truncation
has a unique root $T_K(h,s)=L$ in $[x_0,\infty)$ for every target in
Theorem~\ref{thm:certificate}. This is a globally well-conditioned approximate
inverse on that range.

For an even $K$, monotonicity is guaranteed wherever
$(2K+2)E_K(x)/x<c_{\avec}$. This is an explicit, checkable threshold.
The general residual certificate does not require solving any truncated
equation exactly, and therefore remains useful for Newton iterates, interval
bisection, or a formal inverse series generated by another method.

\subsection{Using formal reversion without confusing it with certification}
Suppose a chosen core $\Phi$ is locally invertible, with $\Phi(s)=L$, and write
$\Delta=F-\Phi$. Classical inverse perturbation gives, formally in an auxiliary
parameter multiplying $\Delta$,
\begin{equation}
 x=s+\sum_{j\ge1}\frac{(-1)^j}{j!}
 \left(\frac1{\Phi'(s)}\frac\dd{\dd s}\right)^{j-1}
                 \frac{\Delta(s)^j}{\Phi'(s)}.
 \label{eq:inverse-formal}
\end{equation}
The repository already develops this kind of reversion calculus
\cite{repo-readme,repo-companion}. The first terms are
\begin{equation}
 x=s-\frac\Delta{\Phi'}
     +\frac{\Delta\Delta'}{(\Phi')^2}
     -\frac{\Delta^2\Phi''}{2(\Phi')^3}+\cdots.
\end{equation}
In the present setting a finite candidate produced from this identity can be
substituted into $T_K$ and certified by \eqref{eq:general-certificate}. We do
not assert convergence of a Bernoulli-based infinite inverse series. The
separation between coefficient production and error certification is precisely
what makes the uniform result usable at the endpoint.

\section{Modular sectors, finite evaluation, and integer decisions}
\label{sec:computation}

\subsection{Two different exponential phenomena}
The modular correction has the convergent expansion
\begin{equation}
 M(h)=\sum_{n\ge1}\sigma_{-1}(n)\e^{-4\pi^2n/h},\qquad
 \sigma_{-1}(n)=\sum_{d\mid n}\frac1d.
 \label{eq:modular-sectors}
\end{equation}
It follows by expanding $-\log(1-\widehat Q^j)$ and grouping terms with the
same product of indices, where $\widehat Q=\e^{-4\pi^2/h}$.
Theorem~\ref{thm:borel} sums the finite-size Bernoulli tail exactly. The full
forward function is that Borel sum plus the separately specified modular
contribution $dM(h)$ and the elementary phase and amplitude.

In the crossover $x=\tau/h$, the uniform least-term envelope has exponential
size $\e^{-2\pi\amin\tau/h}$, whereas the first modular sector has size
$\e^{-4\pi^2/h}$. Their exponents meet at $\tau=2\pi/\amin$. This comparison
is useful for budgeting guaranteed errors; it is not an assertion that an
actual remainder is asymptotic to its upper envelope for all $\tau$. In
particular, keeping the modular exponential while leaving an uncontrolled
algebraic tail does not establish beyond-all-orders accuracy.

\subsection{A finite evaluation budget for the modular factor}
For $0<q<1$,
\begin{equation}
 -\log\prod_{j\ge1}(1-q^j)
          =\sum_{m\ge1}\frac{q^m}{m(1-q^m)}.
\end{equation}
If this sum is cut after $m=J$, its positive remainder is bounded by
\begin{equation}
 \frac{q^{J+1}}{(J+1)(1-q)^2}.
 \label{eq:mod-tail-bound}
\end{equation}
Indeed $m\ge J+1$ and $1-q^m\ge1-q$, after which a geometric sum finishes
the estimate. For small $h$, use $q=\widehat Q$. For large $h$, use
\begin{equation}
 M(h)=-\log\calP(h)+\frac12\log\frac{2\pi}{h}
                  -\frac{\pi^2}{6h}+\frac h{24},
 \label{eq:M-switch}
\end{equation}
so that a product with parameter at least $2\pi$ is always available.
Thus retaining $M$ exactly in the theorem does not require an infinite
computation in a finite certificate.

\subsection{Coefficient evaluation and cancellation control}
For a nonnegative integer $p$,
\begin{equation}
 \Li_{-p}(z)=\frac{zP_p(z)}{(1-z)^{p+1}},
 \label{eq:Eulerian-poly}
\end{equation}
where for $p\ge1$ the coefficients of $P_p$ are Eulerian numbers; $p=0$ is
handled as $z/(1-z)$. The recurrence
\begin{equation}
 \left\langle\genfrac{}{}{0pt}{}{n}{k}\right\rangle
 =(k+1)\left\langle\genfrac{}{}{0pt}{}{n-1}{k}\right\rangle
 +(n-k)\left\langle\genfrac{}{}{0pt}{}{n-1}{k-1}\right\rangle
\end{equation}
gives exact integer coefficient arrays. Thus every correction $D_k$ is
computable by rational operations and real exponentials once the Bernoulli
number is known. Near zero one should evaluate $1-\e^{-s}$ with
\texttt{expm1}, or use the removable-singularity expansion directly.

The remaining dilogarithms can be evaluated by their defining series; for
$0\le z<1$ the tail after $J$ terms is bounded by
\begin{equation}
 \sum_{m>J}\frac{z^m}{m^2}\le
           \frac{z^{J+1}}{(J+1)^2(1-z)}.
\end{equation}
Near $z=1$ this is inefficient. It is preferable to evaluate the balanced phase
via \eqref{eq:Sprime}, its local Taylor series, or standard dilogarithm
functional identities. Fixed working precision without cancellation control
is not a certificate, even when the analytic truncation error is tiny.

\subsection{From a real inverse interval to an integer answer}
When the ray parameters are positive integers, the lattice values at integer
$x$ form an increasing sequence. Let $x_*=F_h^{-1}(\log Y)$. The least index
whose value is at least $Y$ is $\lceil x_*\rceil$. A certified interval
$[u,v]$ determines that index whenever it lies in a single ceiling cell
$(n-1,n]$. If an endpoint touches an integer, the convention at equality must
be respected; exact or validated comparison with that lattice value resolves
the boundary.

This is an application of the repository's general distinction between real
inversion and staircase inversion \cite{repo-readme}, not a replacement for
it. An interval of width less than one alone does not suffice: it may still
straddle an integer threshold.

\section{Verification and reproducibility}
\label{sec:verification}

The accompanying \texttt{verify.py} implements the normalized coefficients,
core, analytic bounds, least-term rule and inverse enclosure formulas using
\texttt{mpmath}. The reference forward values on integral rays are computed
from the finite Gaussian product, independently of the asymptotic expansion;
at $h=0$ the reference is the gamma quotient. This avoids validating an
expansion only against a second evaluation of the same truncated expression.

The recorded run used 240 decimal digits. It checked 812 instances of the
signed forward and first-omitted-term inequalities, including $h=0$,
$h=10^{-8}$, crossover values and $h=10$, for rays $(1,1)$, $(1,2)$,
$(1,1,1)$ and $(1,3,5)$ and indices from 1 to 20. It also checked 24 inverse
enclosures, 36 resonant least-term remainders, 48 peak-sensitive lattice
moments, and auxiliary moment inequalities. Every assertion passed. The
largest observed ratio $|R_K|/E_K$ was approximately $0.92747955$.

\begin{table}[htbp]
\centering\small
\begin{tabular}{rrrrr}
\toprule
$h$ & $x$ & $K$ & $|F-T_K|$ & $E_K(x)$ \\
\midrule
0.1 & 2 & 6 & $7.8486\times10^{-7}$ & $2.0345\times10^{-5}$ \\
0 & 5 & 15 & $3.365\times10^{-15}$ & $2.0289\times10^{-13}$ \\
0.00000001 & 5 & 15 & $3.365\times10^{-15}$ & $2.0289\times10^{-13}$ \\
0.1 & 5 & 15 & $3.365\times10^{-15}$ & $2.0289\times10^{-13}$ \\
2 & 5 & 15 & $3.3487\times10^{-15}$ & $2.0289\times10^{-13}$ \\
10 & 5 & 15 & $6.2724\times10^{-17}$ & $2.0289\times10^{-13}$ \\
0.1 & 10 & 31 & $5.2196\times10^{-29}$ & $6.5431\times10^{-27}$ \\
0.1 & 20 & 62 & $1.9207\times10^{-56}$ & $4.758\times10^{-54}$ \\
\bottomrule
\end{tabular}

\caption{Central-ray least-term checks. The target is the exact finite-product
value at the displayed integral $x$ (the gamma value for $h=0$). Entries are
rounded numerical observations, not outward-rounded interval endpoints.}
\label{tab:numeric}
\end{table}

At $h=0$, $\avec=(1,1)$ and $x=40$, the ratio of the measured absolute error to
the sharp equivalent in \eqref{eq:sharp} is approximately $1.00448044$.
The nearest-odd rule changes discretely with $x$, so the finite-$x$ ratio need
not approach one monotonically. For the fixed resonance $h=2\pi/3$ at
$x=40$, the ratio to \eqref{eq:resonance} is approximately $1.00000514$;
for $h=2\pi$ it agrees with one to the 14 displayed digits.

For a simple inverse example, use the central ray, $h=0$, and target
$L=F_0(5)=\log252$. Solving $T_3(0,s)=L$ gives
$s-5\approx1.1141787\times10^{-8}$; the one-sided analytic enclosure has
width about $1.06666665\times10^{-7}$. For the same index and $h=10^{-4}$,
the observed displacement is $1.1137474\times10^{-8}$. This illustrates the
continuous endpoint transition without changing the certification method.

\paragraph{What these checks establish.}
They are high-precision consistency tests for formulas and signs. They are
not proofs of the theorems, not interval arithmetic, and not Lean kernel
checks. The package's numerical enclosure routine prints high-precision
approximations to analytic enclosure endpoints; those printed endpoints are
not themselves machine-certified until elementary evaluations and rounding
are validated. The mathematical certificate is Theorem~\ref{thm:certificate}.

\paragraph{Reproduction.}
From the package directory run
\begin{lstlisting}
python -m pip install -r requirements.txt
python verify.py
pdflatex -interaction=nonstopmode -halt-on-error uniform_q_multinomial_transseries.tex
pdflatex -interaction=nonstopmode -halt-on-error uniform_q_multinomial_transseries.tex
\end{lstlisting}
The Python run regenerates \texttt{verification\_results.json} and
\texttt{numerical\_table.tex}. The build does not depend on repository-local
notation files or external images. No repository files were modified.

\section{A proposed formalization route}
\label{sec:formalization}

The proof divides naturally into a finite algebraic layer, a real-analysis
layer, and an optional complex-analysis layer. This makes the uniform theorem
a practical next target for the repository's existing residual-certificate
architecture, without requiring a complete formal datatype for transseries.

At the finite layer, formalize the geometric identity that produces
\eqref{eq:kernel-remainder}; the Bernoulli--zeta coefficient identity may be
supplied through a separate analytic theorem. Establish positivity of
$W_{\avec}$ and $-W'_{\avec}$ using only the strict inequalities $a_j<A$.
The moment estimate in Lemma~\ref{lem:moment} is an independent reusable lemma.

At the real-analysis layer, formalize the convergent Bose partial fractions,
the dominated summations in \eqref{eq:Rexact}, and the exact modular identity
as separately named dependencies. Prove Theorem~\ref{thm:uniform} first for
$h>0$, then add the Binet endpoint. A finite truncation object with fields
\texttt{value}, \texttt{errorBound} and \texttt{errorSign} would already support
all inverse certificates. It should not be represented as a convergent
infinite Bernoulli series.

The slope theorem is particularly useful independently of transseries.
Its proof needs decreasing-function Riemann sums and a positive integral
representation of the trigamma combination. Once it is available, the
inverse theorem is a direct mean-value argument; the repository already has
related residual-transport interfaces \cite{repo-readme}.

The Borel pole lattice and the sharp endpoint equivalent are additional
analytic layers. They are not prerequisites for the finite-order inverse
certificate. Separating them avoids inflating a useful elementary theorem's
formal dependency graph. No Lean implementation or successful Lean build is
claimed in this package.

\section{Further research questions and topics}
\label{sec:questions}

The following questions are proposed next steps, not claims that each is an
established open problem in the published literature. Their exact status
would require a focused literature review before being advertised that way.

\begin{enumerate}[leftmargin=*,label=\textbf{Q\arabic*.},itemsep=6pt]
\item \textbf{The lattice-to-integral transition of the least remainder.}
Theorem~\ref{thm:resonance} disproves an $x^{-1/2}$ prefactor uniform in all
$h$. Determine a full transition formula when $h\sqrt{\amin x}$ is bounded,
including the fractional position of the saddle $2\pi/h$ relative to the
integer lattice. A Gaussian lattice sum suggests theta-function factors,
but a uniform theorem must control both the tails and the next partial-fraction
poles. This would interpolate between \eqref{eq:sharp} and the discrete
resonance law rather than incorrectly extrapolating either one.

\item \textbf{A uniform first hyperasymptotic correction.}
Resolve the $\ell=1$, smallest-parameter contribution of
\eqref{eq:Rexact} explicitly, including its dependence on $h$, and then bound
the remainder after subtracting it. The next relevant scales may come from
$2\amin$, the second-smallest ray parameter, and off-axis Borel poles.
A proof should distinguish those mechanisms instead of taking the minimum
of a guessed list of exponents.

\item \textbf{Growing or degenerating rays.}
The theorem fixes $r$ and the positive $a_j$. Determine useful uniform
statements when $r\to\infty$, when $a_*\to0$, or when several parameters
coalesce at an $x$-dependent rate. The explicit sums in $E_K$ expose the
possible loss of constants, but do not by themselves provide an optimal
joint scaling theorem.

\item \textbf{Positive signed product ratios beyond multinomials.}
For a general finite signed exponential weight
$W(s)=\sum_\nu c_\nu\e^{-\lambda_\nu s}$, characterize tractable sufficient
conditions for $W\ge0$ and $-W'\ge0$. Such conditions would transfer the
present remainder and conditioning arguments to wider balanced
$q$-gamma quotients. Necessity of these conditions for alternating remainder
laws is a separate, more delicate question.

\item \textbf{Complex sectors and pole pinching.}
Develop uniform lateral Laplace estimates when $\tau$ approaches zero through
complex directions. The real theorem exploits positivity; it does not cover
contours pinched by moving poles. The meromorphic limit in
Theorem~\ref{thm:nearest} supplies explicit local data for this analysis.

\item \textbf{Complete collision and Stokes-factor classification.}
For commensurable ray parameters, classify all cancellations among the poles
in \eqref{eq:scaled-poles}, including residues at nonzero real translates.
Then prove global Stokes formulas with specified contours and convergence
regions. The nearest pole cannot cancel, but later singularities can.

\item \textbf{Uniform explicit inverse expansions.}
Turn \eqref{eq:inverse-formal} into a single coefficient calculus with
explicit derivative bounds uniform in $h\ge0$. The present article certifies
finite inverse candidates; it does not prove a complete resurgent expansion
of the inverse across every endpoint and complex sector.

\item \textbf{Validated fast computation.}
Implement the phase and coefficient evaluation with directed rounding,
adaptive precision and cancellation-aware formulas. Compare direct products,
$q$-gamma evaluation, normalized asymptotics and hyperasymptotics. The goal
is a fully verified enclosure routine, not merely more digits from an
unvalidated evaluator.

\item \textbf{Summed $q$-multinomials and theta transitions.}
Extend the approach from one multinomial ray to the generalized Galois
numbers and root-lattice theta sectors treated elsewhere in the companion.
Summing over a moving lattice introduces a width transition absent from the
fixed-ray problem. It should not be assumed that the same smallest-parameter
law survives unchanged.

\item \textbf{A formal optimal-truncation certificate.}
Formalize the complete chain from the signed kernel through the explicit
Bernoulli bound and nearest-odd rule to the inverse enclosure. The sharp
endpoint equivalent can be a second-stage target. This would convert the
repository's supporting least-term results into an assembled theorem for a
nontrivial analytic family.
\end{enumerate}

\section{Conclusion}

The explicit repository gap at $hx=0$ can be closed without a new abstract
transseries field. The decisive step is an exact signed kernel identity that
survives the multinomial combination. Normalizing in powers of $1/x$ removes
the artificial endpoint singularities from the coefficient functions and
produces an error bound uniform for every $h\ge0$.

The Borel analysis explains the resulting exponential scale. The nearest
singularities are fixed at $\pm2\pi i a_*$ in the normalized Borel plane;
the gamma endpoint retains them and gives the sharp least-term equivalent.
The peak-sensitive moment bound removes the polynomial loss uniformly, while
fixed lattice resonances prove that a vanishing uniform prefactor is
impossible for the specified truncation rule. The modular Euler-product sectors are separate data and must remain in the
forward model when exponentially small accuracy is requested. Finally, the
uniform derivative comparison turns the forward expansion into a practical
inverse certificate, including the classical endpoint and the fixed-$q$
regime.

These results are a concrete extension of the repository's transseries
program. They do not settle general resurgence or global inversion questions,
and their worldwide novelty is not certified. The strongest claims made here
are the precisely stated theorems with the proofs above.

\appendix
\section{A compact theorem and dependency ledger}

\begin{longtable}{@{}>{\raggedright\arraybackslash}p{.22\textwidth}>{\raggedright\arraybackslash}p{.42\textwidth}>{\raggedright\arraybackslash}p{.28\textwidth}@{}}
\toprule
Result & Content & Main inputs\\\midrule
\endfirsthead
\toprule Result & Content & Main inputs\\\midrule\endhead
Theorem~\ref{thm:uniform} & All-$h$ signed expansion, explicit remainder and derivative bounds & Bose partial fractions, positive weight, moment lemma, eta identity, Binet endpoint\\[4pt]
Theorem~\ref{thm:borel} & Exact single-tail Borel transform and Laplace sum & Geometric series, cosine transform, normally convergent meromorphic sums\\[4pt]
Theorem~\ref{thm:nearest} & Universal nearest poles and endpoint kernel & Explicit Borel formula, residue addition, dominated convergence\\[4pt]
Theorem~\ref{thm:exponential} & Explicit uniform least-term envelope & Bernoulli--zeta identity and real Stirling bound\\[4pt]
Theorem~\ref{thm:constant-envelope} & Uniform $O(\e^{-2\pi a_*x})$ error & Peak-sensitive lattice moment\\[4pt]
Theorem~\ref{thm:resonance} & Sharp fixed-$h$ resonances and the obstruction to a vanishing uniform prefactor & Discrete positive remainder, saddle separation\\[4pt]
Theorem~\ref{thm:sharp} & Sharp endpoint equivalent & Exact Binet remainder, gamma concentration, Stirling formula\\[4pt]
Theorem~\ref{thm:slope} & $F_h'\ge F_0'\ge c_{\avec}$ & Decreasing Riemann sums, digamma and trigamma integrals\\[4pt]
Theorem~\ref{thm:certificate} & One-sided inverse enclosures & Uniform forward error and exact slope bound\\
\bottomrule
\end{longtable}

\section{Conventions and common failure modes}

The expansion variable in Theorem~\ref{thm:uniform} is $1/x$; $hx$ is retained
inside the coefficient functions. The Borel variable $\xi$ initially refers
to $h$, whereas $\eta=\xi/\tau$ refers to $1/x$. Confusing these variables
changes the apparent nearest action.

The approximation $T_K$ contains $K$ Bernoulli terms, ending at $x^{-(2K-1)}$.
The first omitted power is $N=2K+1$. Choosing $N\approx2\pi a_*x$ is therefore
consistent with $K\approx\pi a_*x$, not $K\approx2\pi a_*x$.

The signs are also worth recording. $R_K$ has sign $(-1)^{K+1}$, so odd
truncations lie below the exact forward function and yield upper inverse
approximations. Even truncations lie above it and yield lower inverse
approximations whenever the relevant truncated root is used. The derivative
of $R_K$ has the opposite sign.

The interpolation is fixed by \eqref{eq:exact-F}; a different real
interpolation of the same sequence can have a different inverse between
integer points. All staircase conclusions must use that convention or an
explicit comparison theorem.

Finally, a list of high-precision decimal values is not an interval proof.
A computationally certified inverse needs both the analytic truncation
bound and validated evaluation of the finite expression and its residual.

\clearpage
\phantomsection
\addcontentsline{toc}{section}{References}
\begin{thebibliography}{99}
\bibitem{repo-companion}
Vladimir Reshetnikov, \emph{ProveIt repository: Combinatorial Transseries and
Their Inverses}, source
\texttt{Combinatorial\_Transseries\_Inverses.tex}, inspected at commit
\texttt{04e06e032dff1966513bfba316966d52db3646b4}, 29 September 2026.
Full repository path and inspected labels are recorded in
\texttt{source\_provenance.md}. Repository:
\url{https://github.com/VladimirReshetnikov/ProveIt}.

\bibitem{repo-readme}
Vladimir Reshetnikov, \emph{ProveIt repository: Series and transseries},
\path{Analysis/FabiusFunction/docs/semi-formalized-research-frontiers/}
\path{drafts/series-and-transseries/README.md}, inspected 29 September 2026,
source blob \texttt{516aaa30e84cb0036a00f803cb6e8e3026febb8e}.

\bibitem{costin-garoufalidis}
Ovidiu Costin and Stavros Garoufalidis,
\emph{Resurgence of the Euler--MacLaurin summation formula},
arXiv:math/0703641, version 2 (2007).
\url{https://arxiv.org/abs/math/0703641}.

\bibitem{garoufalidis-kashaev}
Stavros Garoufalidis and Rinat Kashaev,
\emph{Resurgence of Faddeev's quantum dilogarithm},
arXiv:2008.12465 (2020).
\url{https://arxiv.org/abs/2008.12465}.

\bibitem{dlmf-hyperbolic}
NIST Digital Library of Mathematical Functions,
\S4.36, \emph{Infinite Products and Partial Fractions}.
\url{https://dlmf.nist.gov/4.36}.

\bibitem{dlmf-eta}
NIST Digital Library of Mathematical Functions,
\S23.18, \emph{Modular Transformations}.
\url{https://dlmf.nist.gov/23.18}.

\bibitem{dlmf-gamma-int}
NIST Digital Library of Mathematical Functions,
\S5.9, \emph{Integral Representations}.
\url{https://dlmf.nist.gov/5.9}.

\bibitem{dlmf-gamma-asymp}
NIST Digital Library of Mathematical Functions,
\S5.11, \emph{Asymptotic Expansions}.
\url{https://dlmf.nist.gov/5.11}.
\end{thebibliography}
\end{document}
